\begin{abstract}
Exploration via intrinsic rewards can accelerate learning when extrinsic feedback is sparse or deceptive, yet many intrinsic objectives overemphasize novelty in regions of high uncertainty where learning is slow or unstable.
We introduce \emph{Gated Learning-Progress Exploration} (GLPE), a simple intrinsic-reward family that prioritizes transitions that (i) induce meaningful changes in a learned feature representation and (ii) occur in regions where a forward dynamics model is actively improving.
GLPE estimates region-local learning progress using the discrepancy between short- and long-horizon exponential moving averages (EMAs) of prediction error within an online partition of feature space, and combines it with a feature-space impact term.
An optional region-specific gate suppresses intrinsic rewards in regions where prediction error remains high but learning progress is low, providing a targeted guardrail against unproductive curiosity.
We also study \emph{GLPE (no gate)}, which uses the same base intrinsic score without any gating.
We evaluate both variants with a common PPO backbone on five Gymnasium control benchmarks and compare against widely used intrinsic-reward baselines, reporting performance under both step- and wall-clock-normalized views.
Across the suite, GLPE (no gate) matches the strongest baseline within uncertainty, while the full gated variant is most beneficial in sparse-reward regimes where ungated intrinsic rewards can destabilize training.
\end{abstract}
